\chapter{Rotator Cuff Tendinopathy and Impingement}
\index{Rotator Cuff Tendinopathy and Impingement}
\label{sec:Rotator Cuff Tendinopathy and Impingement}

\section{Overview}
Rotator cuff tendinopathy and subacromial impingement syndrome are among the most common causes of shoulder pain, affecting up to 30\% of the population at some point in life. These conditions involve inflammation, degeneration, or mechanical compression of the rotator cuff tendons, particularly the supraspinatus, as they pass through the subacromial space.
\section{Classification}

\begin{description}[font=\normalfont\bfseries,leftmargin=3em,labelwidth=2.5em]
  \item[Cause] Degenerative / Mechanical
  \item[Organ System] Bones / Musculoskeletal
  \item[Pathophysiology] Tendon degeneration, inflammation, and mechanical compression within the subacromial space
  \item[Duration] Chronic (months to years)
  \item[Onset] Gradual (or acute if superimposed tear)
  \item[Mode of Transmission] Non-communicable
  \item[Clinical Course] Chronic, progressive without treatment
  \item[Severity] Mild to Severe
  \item[Extent of Disease] Localized
  \item[Age Group] Adults (40--60 years)
  \item[Epidemiology] Very common; 30--50\% of all shoulder pain presentations
  \item[ICD-11 Code] FA52.0
\end{description}

\section{Causes}
Rotator cuff tendinopathy results from a combination of intrinsic tendon degeneration (age-related, hypovascularity) and extrinsic mechanical compression from the acromion, coracoacromial ligament, or acromioclavicular joint. Repetitive overhead activities, tendon overload, and poor scapular mechanics contribute to progressive tendon damage. Partial or full-thickness tears may develop as the condition advances.

\section{Risk Factors}

\begin{itemize}[noitemsep]
  \item Age over 40 years
  \item Repetitive overhead activities (painting, swimming, throwing sports)
  \item Occupational overhead work (construction, electricians)
  \item Poor posture and scapular dyskinesis
  \item Acromial morphology (curved or hooked acromion)
  \item Smoking
  \item Diabetes mellitus
  \item Previous shoulder injury
\end{itemize}

\section{Signs and Symptoms}

\subsection{Early symptoms}
\begin{itemize}[noitemsep]
  \item Early manifestations of rotator cuff tendinopathy and impingement may be subtle and nonspecific
\end{itemize}

\subsection{Common symptoms}
\begin{itemize}[noitemsep]
  \item \subsection{Early symptoms}
  \item \subsection{Common symptoms}
  \item \textbf{Common symptoms:}
  \item Deep, aching shoulder pain, especially with overhead activities
  \item Night pain that disrupts sleep, particularly when lying on the affected side
  \item Pain with arm elevation between 60 and 120 degrees (painful arc sign)
  \item Weakness in shoulder abduction and external rotation
  \item Positive impingement signs (Neer, Hawkins-Kennedy)
  \item Stiffness and reduced range of motion
  \item Pain or discomfort in the affected area
  \item Difficulty performing daily activities
\end{itemize}

\subsection{Less common symptoms}
\begin{itemize}[noitemsep]
  \item Atypical or less common presentations of rotator cuff tendinopathy and impingement may occur
\end{itemize}

\subsection{Emergency warning signs}
\begin{itemize}[noitemsep]
  \item Severe or worsening symptoms of rotator cuff tendinopathy and impingement require immediate medical evaluation
\end{itemize}
\section{Progression and Stages}
The condition progresses through three stages: Stage 1 involves edema and hemorrhage (reversible); Stage 2 involves tendinosis and fibrosis (chronic); Stage 3 involves partial or full-thickness tendon tears. Without intervention, progressive tendon degeneration leads to worsening pain, weakness, and functional limitation.

\section{Complications}

\begin{itemize}[noitemsep]
  \item Progression to full-thickness rotator cuff tear
  \item Muscle atrophy and fatty infiltration of rotator cuff muscles
  \item Frozen shoulder (adhesive capsulitis) from disuse
  \item Chronic pain and disability
  \item Reduced quality of life
  \item Emotional and psychological effects
\end{itemize}

\section{Diagnosis}

\begin{itemize}[noitemsep]
  \item Clinical diagnosis based on history and physical examination
  \item Impingement tests (Neer, Hawkins-Kennedy, painful arc)
  \item Rotator cuff strength testing (Jobe test, external rotation lag sign)
  \item Ultrasound or MRI to assess tendon integrity and identify tears
  \item X-ray to evaluate acromial morphology and exclude other pathology
\end{itemize}

\section{Prevention}

\begin{itemize}[noitemsep]
  \item Maintain good shoulder posture and ergonomics
  \item Strengthen rotator cuff and scapular stabilizer muscles
  \item Avoid sudden increases in overhead activity intensity
  \item Use proper technique in sports and occupational activities
  \item Go for regular check-ups so any problems are caught early
\end{itemize}

\section{Treatment Options}

\begin{itemize}[noitemsep]
  \item Physical therapy focusing on rotator cuff strengthening and scapular stabilization
  \item Activity modification and avoidance of overhead aggravating activities
  \item Nonsteroidal anti-inflammatory drugs (NSAIDs) for pain and inflammation
  \item Subacromial corticosteroid injections for symptomatic relief
  \item Platelet-rich plasma (PRP) injections (emerging evidence)
  \item Arthroscopic subacromial decompression for refractory cases
  \item Rotator cuff repair surgery for full-thickness tears
  \item Lifestyle changes and self-care
\end{itemize}

\section{Living with It}

\begin{itemize}[noitemsep]
  \item Follow your treatment plan as prescribed
  \item Keep up with regular appointments with your healthcare provider
  \item Adjust your daily habits as needed
  \item Reach out to family, friends, or support groups for help
  \item Keep learning about your condition and how to manage it
\end{itemize}

\section{Outlook}
Most patients with rotator cuff tendinopathy improve with conservative management within 6--12 weeks. Outcomes are better with early intervention. Full-thickness tears may require surgical repair to restore function. Chronic or massive tears can lead to irreversible muscle atrophy and persistent disability, emphasizing the importance of timely diagnosis and treatment.
